%----------------------------------------------------------
% Yash Patel — Resume (modern minimal)
% Compile with pdflatex
%----------------------------------------------------------
\documentclass[letterpaper,12pt]{article}

\usepackage[top=0.39in,bottom=0.39in,left=0.47in,right=0.47in]{geometry}
\usepackage[T1]{fontenc}
\usepackage[utf8]{inputenc}
\usepackage{charter}  % clean serif (Bitstream Charter)
\usepackage{microtype}
\usepackage{enumitem}
\usepackage{titlesec}
\usepackage{xcolor}
\usepackage[hidelinks]{hyperref}
\usepackage{tabularx}
\input{glyphtounicode}
\pdfgentounicode=1

\definecolor{accent}{HTML}{1A1A1A}
\definecolor{rulecolor}{HTML}{C8C8C8}
\definecolor{muted}{HTML}{555555}

\pagestyle{empty}
\setlength{\parindent}{0pt}
\raggedright
\raggedbottom

% Section: small caps, letterspaced, thin gray rule
\titleformat{\section}
  {\scshape\large\color{accent}}
  {}{0em}{}
  [{\color{rulecolor}\titlerule[0.5pt]}]
\titlespacing*{\section}{0pt}{2pt}{1pt}

% Entry heading: bold left, dates right; role/location line below
\newcommand{\entry}[4]{%
  \begin{tabularx}{\textwidth}{@{}X@{}r@{}}
    \textbf{#1} & {\small\color{muted}#2} \\
    {\small\itshape #3} & {\small\itshape\color{muted}#4} \\
  \end{tabularx}\vspace{3pt}
}

\newcommand{\project}[3]{%
  \begin{tabularx}{\textwidth}{@{}X@{}r@{}}
    \textbf{#1}\; {\small\color{muted}\itshape #2} & {\small #3} \\
  \end{tabularx}\vspace{3.5pt}
}

\newenvironment{bullets}
  {\begin{itemize}[leftmargin=1.1em, itemsep=0pt, topsep=0pt, parsep=0pt, label={\color{muted}\footnotesize\textbullet}]\fontsize{11}{13.0}\selectfont}
  {\end{itemize}\vspace{-3pt}}

\newcommand{\skill}[2]{\textbf{#1}\enspace {\color{muted}·}\enspace #2\\[-1pt]}

\begin{document}

%---------- Header (centered) ----------
\begin{center}
  {\huge\scshape Yash Patel}\\[1pt]
  {\small
    +16467735448\,{\color{muted}·}\,
    \href{mailto:yhp2378@gmail.com}{yhp2378@gmail.com}
    \,{\color{muted}·}\,
    \href{https://www.linkedin.com/in/yashp5/}{linkedin.com/in/yashp5}
    \,{\color{muted}·}\,
    \href{https://github.com/yashp5}{github.com/yashp5}
  }\\[2pt]
  {\fontsize{9}{10.5}\selectfont\color{muted}
    \textbf{Work authorization:}\enspace
    \ifdefined\CANADA
      Canada --- employer-sponsored permit required; GTS Category B (NOC 21231/21232) may apply.
    \else
      UK --- eligible to apply for HPI (NYU M.S., 2025); no employer sponsorship for HPI; visa not yet held.
    \fi
  }
\end{center}
\vspace{-12pt}

%---------- Skills ----------
\section{Skills}\vspace{1.2pt}
{\small
\skill{Languages}{Go, Python, C++, TypeScript}
\skill{LLM Serving}{llama.cpp, Continuous Batching, KV-Cache Management, Inference Scheduling}
\skill{Distributed Systems}{gRPC, Concurrency, WebSockets, Kafka, HTTP APIs}
\skill{Infrastructure}{Docker, Kubernetes, Linux, AWS, GCP}
\skill{Data \& Observability}{PostgreSQL, Redis, Prometheus, Grafana}
}

%---------- Experience ----------
\section{Experience}

\entry{Dora}{Jun 2025 -- Feb 2026}{Software Engineer}{New York}
\begin{bullets}
  \item Built low-latency Go services for a tokenized bond exchange, including order matching, market data, and account updates.
  \item Reduced matching-loop p50 latency 19\% (4.55ms to 3.69ms) by using event-driven in-memory caches instead of database lookups.
  \item Built concurrent WebSocket streaming APIs with shared caches and fan-out; used PostgreSQL LISTEN/NOTIFY for event-driven updates.
\end{bullets}

\vspace{7pt}
\entry{Cashflo}{Sep 2022 -- Sep 2023}{Software Engineer}{Mumbai}
\begin{bullets}
  \item Built a distributed Go ingestion pipeline with batched uploads, sequence tracking, idempotent retries, and cross-instance deduplication.
  \item Designed ERP record reconciliation with JSON schema validation and Kafka routing for valid and invalid records.
\end{bullets}

\vspace{7pt}
\entry{Elocity}{May 2020 -- Sep 2022}{Software Engineer}{Bangalore}
\begin{bullets}
  \item Built a platform supporting 700+ persistent WebSocket connections from EV chargers, with Kafka telemetry and session aggregation.
  \item Supported Docker, Helm, AWS/Kubernetes, and GitLab CI deployments, with Istio and Grafana observability.
\end{bullets}

%---------- Projects ----------
\section{Projects}

\project{Inference Engine}{Go, Python, gRPC}{\href{https://github.com/yashp5/inference-engine}{github.com/yashp5/inference-engine}}
\begin{bullets}
  \item Built an LLM serving engine with a Go control plane and gRPC Python llama.cpp workers for request scheduling, rate limiting, and admission control.
  \item Implemented continuous batching with per-sequence KV-cache lifecycle management; improved output throughput 49\% (94 to 140 tok/s) and cut p99 end-to-end latency 40\% (55.4s to 33.3s) versus static admission on a mixed-length workload at concurrency 32.
  \item Built Go end-to-end and Python engine/gRPC tests and a benchmark suite measuring throughput, time-to-first-token, latency, and slot occupancy.
\end{bullets}

\vspace{6pt}
\project{GPU Demand Response Protocol}{Go, WebSockets}{\href{https://github.com/grene-labs/grene}{github.com/grene-labs/grene}}
\begin{bullets}
  \item Built a GPU power-curtailment control plane with agent telemetry, heartbeats, automated cap restoration, and hardware-free tests.
\end{bullets}

%---------- Education ----------
\section{Education}
\textbf{New York University}\hspace{.5em} {\small\itshape M.S. in Computer Engineering}\hfill{\small\color{muted}2023 -- 2025}\\[1pt]
\textbf{Ahmedabad University}\hspace{.5em} {\small\itshape B.S. in Information and Communication Technology}\hfill{\small\color{muted}2016 -- 2020}

\end{document}
